% SEGMENT OLP-0092-S01
% Part: first-order-logic
% Chapter: natural-deduction
% Section: provability-consistency

\documentclass[../../../include/open-logic-section]{subfiles}

\begin{document}

\iftag{FOL}
      {\olfileid{fol}{ntd}{prv}}
      {\olfileid{pl}{ntd}{prv}}
      
\olsection{\usetoken{S}{derivability} மற்றும் முரணின்மை}

!!{derivability} உறவின் பல பண்புகளை இப்போது நிறுவுவோம். அவை தனித்தும்
சுவையானவை; மேலும் அவை ஒவ்வொன்றும் நிறைவுத் தேற்றத்தின் நிறுவலில் ஒரு பங்கு
வகிக்கும்.

\begin{prop}\ollabel{prop:provability-contr}
  $\Gamma \Proves !A$ என்றும், $\Gamma \cup \{!A\}$ முரணுடையது என்றும்
  இருந்தால், $\Gamma$ முரணுடையது.
\end{prop}

\begin{proof}
$\Gamma$ இலிருந்து~$!A$ இன் !!{derivation} ஐ~$\delta_1$ என்றும்,
$\Gamma \cup \{!A\}$ இலிருந்து~$\lfalse$ இன் !!{derivation} ஐ~$\delta_2$
என்றும் கொள்க. அப்போது பின்வருமாறு !!{derive} இயலும்:
\begin{prooftree}
\AxiomC{$\Gamma, \Discharge{!A}{1}$}
\RightLabel{$\delta_2$}
\DeduceC{$\lfalse$}
\DischargeRule{\Intro{\lnot}}{1}
\UnaryInfC{$\lnot !A$}
\AxiomC{$\Gamma$}
\RightLabel{$\delta_1$}
\DeduceC{$!A$}
\RightLabel{\Elim{\lnot}}
\BinaryInfC{$\lfalse$}
\end{prooftree}
புதிய !!{derivation}[loc] எடுகோள்~$!A$ ஆனது !!{discharged}; ஆகவே இது
$\Gamma$ இலிருந்து !!a{derivation}.
\end{proof}

% SEGMENT OLP-0092-S02
\begin{prop}
\ollabel{prop:prov-incons}
$\Gamma \Proves !A$ என்பது, அப்போதும் அப்போது மட்டுமே, $\Gamma \cup
\{\lnot !A\}$ முரணுடையது.
\end{prop}

\begin{proof}
முதலில் $\Gamma \Proves !A$ என்க; அதாவது, !!{undischarged}
எடுகோள்கள்~$\Gamma$ இலிருந்து~$!A$ இன் !!a{derivation}~$\delta_0$ உள்ளது.
$\Gamma \cup \{\lnot !A\}$ இலிருந்து~$\lfalse$ இன் !!a{derivation} ஒன்றைப்
பின்வருமாறு பெறுகிறோம்:
\begin{prooftree}
  \AxiomC{$\lnot !A$}
  \AxiomC{$\Gamma$}
  \RightLabel{$\delta_0$}
  \DeduceC{$!A$}
  \RightLabel{\Elim{\lnot}}
  \BinaryInfC{$\lfalse$}
\end{prooftree}

இப்போது $\Gamma \cup \{\lnot !A\}$ முரணுடையது என்க; அதற்குரிய
$\Gamma \cup \{\lnot !A\}$ விலுள்ள !!{undischarged} எடுகோள்களிலிருந்து
$\lfalse$ இன் !!{derivation} ஐ~$\delta_1$ என்க. $\FalseCl$ ஐப் பயன்படுத்தி
$\Gamma$ விலிருந்து மட்டும்~$!A$ இன் !!a{derivation} ஒன்றைப் பெறுகிறோம்:
\begin{prooftree}
  \AxiomC{$\Gamma, \Discharge{\lnot !A}{1}$}
  \RightLabel{$\delta_1$}
  \DeduceC{$\lfalse$}
  \RightLabel{\FalseCl}
  \DischargeRule{\FalseCl}{1}
  \UnaryInfC{$!A$}
\end{prooftree}
\end{proof}

% SEGMENT OLP-0092-S03
\begin{prob}
$\Gamma \Proves \lnot !A$ என்பது, அப்போதும் அப்போது மட்டுமே, $\Gamma
\cup \{!A\}$ முரணுடையது என்பதை நிறுவுக.
\end{prob}

% SEGMENT OLP-0092-S04
\begin{prop}\ollabel{prop:explicit-inc}
  $\Gamma \Proves !A$ என்றும் $\lnot !A \in \Gamma$ என்றும் இருந்தால்,
  $\Gamma$ முரணுடையது.
\end{prop}

\begin{proof}
  $\Gamma \Proves !A$ என்றும் $\lnot !A \in \Gamma$ என்றும் கொள்க.
  $\Gamma$ இலிருந்து~$!A$ இன் !!a{derivation}~$\delta$ உள்ளது.
  $\Elim{\lnot}$ விதியின் இந்த எளிய பயன்பாட்டைக் கருதுக:
  \begin{prooftree}
    \AxiomC{$\lnot !A$}
    \AxiomC{$\Gamma$}
    \RightLabel{$\delta$}
    \DeduceC{$!A$}
    \RightLabel{\Elim{\lnot}}
    \BinaryInfC{$\lfalse$}
  \end{prooftree}
  $\lnot !A \in \Gamma$ என்பதால், அனைத்து !!{undischarged} எடுகோள்களும்
  $\Gamma$ வில் உள்ளன; ஆகவே இது $\Gamma \Proves \lfalse$ என்பதைக் காட்டுகிறது.
\end{proof}

% SEGMENT OLP-0092-S05
\begin{prop}\ollabel{prop:provability-exhaustive}
  $\Gamma \cup \{!A\}$, $\Gamma \cup \{\lnot !A\}$ ஆகிய இரண்டும்
  முரணுடையவை எனில், $\Gamma$ முரணுடையது.
\end{prop}

\begin{proof}
$\Gamma \cup \{ !A \}$ இலிருந்து~$\lfalse$ உம், $\Gamma \cup \{ \lnot !A
\}$ இலிருந்து~$\lfalse$ உம் ஆகியவற்றின் !!{derivation}s முறையே~$\delta_1$,
$\delta_2$ உள்ளன. அப்போது பின்வருமாறு !!{derive} இயலும்:
\begin{prooftree}
\AxiomC{$\Gamma, \Discharge{\lnot !A}{2}$}
\RightLabel{$\delta_2$}
\DeduceC{$\lfalse$}
\DischargeRule{\Intro{\lnot}}{2}
\UnaryInfC{$\lnot \lnot !A$}
\AxiomC{$\Gamma, \Discharge{!A}{1}$}
\RightLabel{$\delta_1$}
\DeduceC{$\lfalse$}
\DischargeRule{\Intro{\lnot}}{1}
\UnaryInfC{$\lnot !A$}
\RightLabel{\Elim{\lnot}}
\BinaryInfC{$\lfalse$}
\end{prooftree}
எடுகோள்கள்~$!A$, $\lnot !A$ ஆகியவை !!{discharged} ஆவதால், இது
$\Gamma$ இலிருந்து மட்டும்~$\lfalse$ இன் !!a{derivation}. ஆகவே $\Gamma$
முரணுடையது.
\end{proof}

\end{document}
